\chapter{Set-up and Measurement}\label{r03:ch:setup}

\emph{The test system, the benchmark application and its variants, the kernel lab and the
probes: how every number in the rest of the report was measured.}

\section{The test system}\label{r03:sec:hw}

Two computers were used. The \term{board} runs every speed measurement. The
\term{laptop}, an x86-64 machine, generates code, cross-compiles, runs QEMU and runs the
waLBerla test suite.

\begin{table}[htbp]
  \centering\small
  \caption{The test system. The board facts were read on the board itself
    (\texttt{/proc/cpuinfo}, the frequency governor, the compiler versions).}\label{r03:tab:hw}
  \begin{tabularx}{\linewidth}{@{}lX@{}}
    \toprule
    Board & Banana~Pi BPI-F3 with a SpacemiT~K1: 8 SpacemiT~X60 cores
      (\code{uarch: spacemit,x60}) in two clusters of 4\\
    Core & in-order RISC-V core with RVV~1.0, VLEN~256; 1.6~GHz, frequency fixed by the
      \code{performance} governor\\
    Caches & L1D 32~KiB per core (4-way set associative); L2 512~KiB per cluster
      (16-way); no L3\\
    Memory & 4~KiB pages; transparent huge pages in \code{madvise} mode\\
    ISA string & lists \code{zicbom} and \code{zicboz}, \emph{not} \code{zicbop}
      (\cref{r03:sec:bg-rvv})\\
    Board compilers & GCC 14.2.0 and Clang 18.1.8, as shipped with the board's Linux
      distribution (Bianbu)\\
    Laptop tools & GCC 16.2.1 cross compiler for RISC-V (\code{riscv64-linux-gnu-g++}), used to
      read the generated assembly; QEMU~11.1.3 (\code{qemu-riscv64 -cpu
      rv64,v=true,vlen=256,zicbop=true}); an x86 waLBerla build with MPI for the tests\\
    Profiling & \code{perf} (\cref{r03:sec:bg-measure}); hardware-event sampling gives no
      samples on this board, so profiles sample a timer instead (\code{-e cpu-clock}); the
      X60's vendor events are counted with \ghws{misc/x60-lab/pmu_stat.sh}\\
    \bottomrule
  \end{tabularx}
\end{table}

\paragraph{How jobs run on the board.} The board is shared with other users. Every
measurement therefore runs as an unattended shell script, called a \term{session script}
(\file{misc/x60-lab/queue_*.sh}), which:
\begin{itemize}
  \item works only in one user directory and writes its logs there;
  \item pins the program to one core, core~2 (\code{taskset -c 2}), so that the operating
        system does not move it between cores;
  \item runs alone: it starts in a detached \code{screen} terminal session and first waits
        until the previous session script has finished, so that two measurements never
        compete for the memory system;
  \item keeps running when the network connection to the board drops.
\end{itemize}
One run of a session script is a \term{session} (\cref{r03:sec:bg-measure}).

\section{The benchmark application}\label{r03:sec:app}

\ghwl{apps/benchmarks/SweepgenLbmCPU/SweepgenLbmCPU.cpp} is a waLBerla benchmark built for
this purpose:
\begin{itemize}
  \item D3Q19, SRT collision, compressible equilibrium, with the physical parameters of
        waLBerla's \code{DoubleShearLayer} example: Reynolds number $30\,000$, initial
        velocity amplitude $u_0 = 0.04$ (a shear-layer flow, so the velocity field is not
        trivial);
  \item a fully periodic domain with one block of $N^3$ cells per MPI process;
  \item arguments: cells per block and direction (default 64), timed time steps (200) and
        untimed warm-up steps (10).
\end{itemize}
At the end it prints one line with the MLUP/s and a \term{checksum}: the sum of $|u|^2$ over
all cells, which the program calls the kinetic energy. The sum is formed in double precision
even in single-precision builds. Two variants that compute the same thing must print the
same checksum, up to rounding.

\paragraph{Variants.} Each way of generating the kernel is a separate executable, defined in
\ghwl{apps/benchmarks/SweepgenLbmCPU/CMakeLists.txt}. The executable names are
\code{SweepgenLbmCPU_} followed by a variant name, which is built from these parts:
\begin{description}
  \item[\code{scalar}] plain C++ kernel, compiled with \code{-march=rv64gc}: no vector
        instructions at all.
  \item[\code{autovec}] plain C++ kernel, compiled with waLBerla's flags, which allow the
        compiler to auto-vectorize.
  \item[\code{explicit}] RVV kernel: \ps{} writes the vector intrinsics
        (\cref{r03:sec:ps-simd}).
  \item[\code{_cse}] lbmpy's global common-subexpression elimination is on.
  \item[\code{_aa}] AA streaming instead of pull (\cref{r03:sec:bg-stream}).
  \item[\code{_pf}] software prefetching 128~bytes ahead (\cref{r03:sec:ps-prefetch}).
  \item[\code{_split}] the direction-split AA kernel (\cref{r03:sec:wl-split}).
\end{description}

\begin{table}[htbp]
  \centering\small
  \caption{The variants of the benchmark at \texttt{efa2b70c}, and the names used for them in
    the text.}\label{r03:tab:variants}
  \begin{tabularx}{\linewidth}{@{}lX@{}}
    \toprule
    \textbf{Variant} & \textbf{Kernel, and name in the text}\\
    \midrule
    \code{scalar} & plain C++, pull: the \textbf{plain scalar build}, the baseline of the
      goal\\
    \code{scalar_cse} & plain C++ with CSE, pull\\
    \code{scalar_cse_aa} & plain C++ with CSE and AA streaming: the \textbf{best scalar
      build}, which has every non-vector optimisation the application offers\\
    \code{autovec}, \code{autovec_cse} & plain C++, compiler may auto-vectorize\\
    \code{explicit}, \code{explicit_cse} & RVV, pull: the vector kernel as it was before this
      work\\
    \code{explicit_pf}, \code{explicit_cse_pf} & RVV, pull, software prefetch (``pull +
      prefetch'')\\
    \code{explicit_cse_aa} & RVV, AA streaming, no prefetch\\
    \code{explicit_cse_aa_pf} & RVV, AA streaming, software prefetch (``AA + prefetch''): the
      fastest variant\\
    \code{explicit_cse_aa_split} & RVV, AA streaming with the direction-split kernel\\
    \bottomrule
  \end{tabularx}
\end{table}

\Cref{r03:fig:timestep} (\cref{r03:sec:bg-stream}) shows the timed loop of the pull and the
AA variants.

\paragraph{Protocol.} \ghws{scripts/walberla_lbm_bench.sh}\code{ 64 200 3} configures one
build directory, builds the variants, runs each 3 times on core~2 and prints the best MLUP/s
and the checksum of each. Unless a table says otherwise, every end-to-end number in this
report is such a best-of-3 value at $64^3$ cells, 200 timed steps after 10 warm-up steps.
\term{End to end} means the whole application: the kernel, the halo exchange and the time
loop.

\section{The kernel lab}\label{r03:sec:lab}

A waLBerla build with all variants takes hours on the board. The \term{kernel lab}
(\ghws{misc/x60-lab/lab/}) instead times one kernel in a small driver program, so that a
change can be tested in minutes. It has these parts:
\begin{description}
  \item[\ghws{misc/x60-lab/lab/gen.py}] builds the update exactly as sweepgen's \code{LbmBulk}
        does for the benchmark, and lets \ps{} generate the kernel. The generated loop bodies
        are identical to those in the waLBerla build. These are the \term{generated kernels}.
  \item[\file{misc/x60-lab/lab/hand/}] \term{hand kernels}: D3Q19 kernels written by hand
        with RVV intrinsics, to try code structures that \ps{} cannot generate.
  \item[\ghws{misc/x60-lab/lab/driver.cpp}] lays out the fields as waLBerla does on RISC-V:
        fzyx layout, one ghost layer, rows padded to 16~bytes (68 floats for 64 cells plus 2
        ghost cells), interior aligned to 16~bytes. Between time steps it refreshes the ghost
        layers periodically; this part is not timed. It counts cycles and instructions around
        the kernel calls only, and prints three checksums: the sum of $|u|^2$, the sum of the
        density field (\term{mass}), and a \term{weighted PDF sum}: the stored PDF values of
        all cells, each multiplied by its direction number plus one, $\sum_{\text{cells}}
        \sum_{i=0}^{18} (i+1)\,f_i$. The last one changes when any single PDF changes, even
        if density and velocity do not. Environment variables change the layout (\code{ALIGN},
        \code{XPAD}, \code{PADF}, \code{THP}); \code{AA=1} runs AA kernels.
  \item[\ghws{misc/x60-lab/lab/edits.py}] applies source-level changes to a generated kernel,
        so that an optimisation can be tried before it is implemented in \ps{}.
  \item[\ghws{misc/x60-lab/lab/run.sh}] builds and runs kernels on the board and appends one
        line per run to \file{results.tsv}.
\end{description}
Each lab kernel has a file name that encodes how it was made; \cref{r03:sec:repro-gen}
explains the naming.

\paragraph{Acceptance test.} Before it was used for any decision, the lab had to reproduce
waLBerla's speed (f32, GCC). The RVV kernel with CSE took 653 cycles per cell in the lab and
682 in waLBerla, which also does the halo exchange; the scalar kernel with CSE took 1035 and
1033. The values agree within 5\,\%, so lab results carry over to the application.

\begin{pitfall}[For AA kernels, the fastest step overstates the speed]
The driver first reported the fastest and the median time step. An AA kernel alternates an
even and an odd step, which are two different functions with different costs, and the
fastest step is always an even one. The fair rate is the \term{cycle average}: 1.6~GHz divided
by the cycles per cell counted over all steps. The driver now also prints the rate over all
steps (\code{mean_mlups}): the same quantity, measured with the clock instead of the cycle
counter.
For example, the generated f32 AA kernel with prefetch: fastest step 8.99, cycle average
8.33~MLUP/s (\cref{r03:tab:labapp}). Tables in this report say which of the two they show.
\end{pitfall}

\section{Probes}\label{r03:sec:probes}

Two small programs (\file{misc/x60-lab/probes/}) measure the core itself, independent of any
LBM code:
\begin{itemize}
  \item \ghws{misc/x60-lab/probes/mem.cpp} measures the DRAM latency, how many misses
        overlap, the throughput of many concurrent streams, and the effect of prefetch and
        cache-block instructions;
  \item \ghws{misc/x60-lab/probes/ops.cpp} measures the throughput and latency of single
        vector and scalar instructions.
\end{itemize}
Both read the cycle counter through the Linux interface \code{perf_event_open}
(\ghws{misc/x60-lab/probes/pmu.hpp}).
